% Soluzioni degli esercizi del Modulo 9
\chapter{Soluzioni -- Modulo 9 -- Ereditarietà, Polimorfismo, Classi Astratte e RTTI}

\textbf{Esercizio 9.1} --- Costruttori dalla base alla derivata, distruttori al
contrario, alla chiusura del blocco.
\codice{codice/m09/es01.cpp}
\uscita{codice/m09/output/es01.txt}

\textbf{Esercizio 9.2} --- Le righe commentate sono quelle che non compilano.
\codice{codice/m09/es02.cpp}
\uscita{codice/m09/output/es02.txt}

\textbf{Esercizio 9.3} --- Tramite \cc|Veicolo*| il metodo non virtuale esegue la
versione di \cc|Veicolo|, quello virtuale la versione di \cc|Moto|. Sulla copia
\cc|Veicolo copia = m;| anche il virtuale esegue la versione di \cc|Veicolo|: la copia
\emph{è} un \cc|Veicolo| (slicing).
\codice{codice/m09/es03.cpp}
\uscita{codice/m09/output/es03.txt}

\textbf{Esercizio 9.4}
\codice{codice/m09/es04.cpp}
\uscita{codice/m09/output/es04.txt}

\textbf{Esercizio 9.5} --- \cc|Color(c)| converte l'intero letto in un valore
dell'\cc|enum|. L'array di puntatori è dinamico perché la sua dimensione si conosce solo
leggendo l'input: si libera prima ogni figura, poi l'array.
\codice{codice/m09/es05.cpp}
\ingresso{codice/m09/input/es05.txt}
\uscita{codice/m09/output/es05.txt}

\textbf{Esercizio 9.6} --- Senza \cc|virtual| il distruttore della derivata non parte e
il suo buffer non viene liberato.
\codice{codice/m09/es06.cpp}
\uscita{codice/m09/output/es06.txt}

\textbf{Esercizio 9.7} --- \cc|typeid(...).name()| restituisce \texttt{4Auto} con g++:
un nome interno del compilatore, utile solo per il debug.
\codice{codice/m09/es07.cpp}
\uscita{codice/m09/output/es07.txt}

\textbf{Esercizio 9.8} --- \cc|to_string| converte il numero in testo; \cc|c_str()|
restituisce la stringa nel formato \cc|const char*| che \cc|what()| richiede. Il messaggio
è salvato in un attributo, così resta valido anche dopo la fine del costruttore.
\codice{codice/m09/es08.cpp}
\uscita{codice/m09/output/es08.txt}

\textbf{Esercizio 9.9} --- Si ordinano i \emph{puntatori}: gli oggetti restano dove
sono nello heap, cambia solo l'ordine delle frecce.
\codice{codice/m09/es09.cpp}
\uscita{codice/m09/output/es09.txt}

\textbf{Esercizio 9.10} --- Senza \cc|virtual| il costruttore di \cc|A| parte due volte:
in \cc|D| ci sono due \cc|A|. Con l'ereditarietà virtuale una volta sola.
\codice{codice/m09/es10.cpp}
\uscita{codice/m09/output/es10.txt}
